% s5_discussion.tex
% Status: DRAFT — generated by TAR-Author; review before locking
% Lock condition: after Phase 15 class-incremental results available

\section{Discussion}
\label{sec:discussion}

\subsection{Limitations}

\textbf{Benchmark scope.}
All experiments use Split-CIFAR-10 in the task-incremental setting with task identity
provided at test time.
This is the easiest of the three continual learning scenarios~\cite{vandeVen2019}
and may overstate absolute performance relative to class-incremental or
domain-incremental evaluations.
TCL's performance in harder settings has not been evaluated and is the subject of
ongoing investigation (Phase~15).

\textbf{Dataset scale.}
Results on a single benchmark cannot confirm generality.
Split-CIFAR-100 (20 five-way tasks) and Split-TinyImageNet would test TCL under
longer task sequences and higher inter-task similarity; both are left to future work.

\textbf{Statistical power.}
With $n = 5$ seeds, the comparison between full TCL and penalty-only does not reach
significance ($p = 0.315$, $d = 0.51$).
Resolving this comparison requires more seeds before a mechanistic claim about the
incremental value of the full loop can be justified.
The pre-registered comparisons against SGD and EWC are both significant.

\textbf{Mechanism.}
The ablation indicates that the anchor penalty is load-bearing and that the
governor-only path is insufficient, but the causal mechanism has not been
confirmed experimentally.
The current five-seed data do not support a simple variance-compression story:
full TCL std is $0.022$ versus $0.018$ for penalty-only. What the ablation does
show is that full TCL attains the lowest mean forgetting, while the marginal gain
over penalty-only remains modest at the present sample size.

\textbf{D$_{\mathrm{PR}}$ disabled.}
The participation ratio scaling was disabled in these experiments for computational
efficiency.
Whether online $D_{\mathrm{PR}}$ computation provides meaningful benefit over the
$D_{\mathrm{PR}} = 1.0$ baseline is not tested here.

\subsection{The SI Collapse as a Methodological Finding}

The SI result (0.500 accuracy on 5/5 seeds) is not a criticism of SI as an algorithm.
It demonstrates a failure mode of published-default benchmarking: hyperparameters that
perform well in one setting can produce degenerate solutions in another.
At $c = 0.1$, $\xi = 0.001$, importance weights accumulate to the point where the
network cannot update meaningfully after the first task.
The network does not forget because it never learns.

The forgetting metric --- the standard ranking criterion in continual learning ---
is blind to this failure: a constant-output network has zero forgetting.
The JAF metric we propose is a minimal corrective that penalises accuracy collapse
without introducing new hyperparameters.
We do not claim JAF is the only solution; we claim it is the minimum necessary
extension to forgetting-only evaluation.

Our SI robustness sweep (Appendix~\ref{app:si_sweep}) found that $c = 0.01$
recovers non-trivial accuracy, whereas $c \in \{0.1, 0.5\}$ remains collapsed.
The failure is therefore specific to part of the published default neighbourhood
rather than to the entire SI family.

\subsection{Future Directions}

\textbf{Fisher-weighted anchor penalty.}
Replacing $D_{\mathrm{PR}}$ with a diagonal Fisher estimate would bring the anchor
penalty closer to EWC in its importance-weighting mechanism, potentially capturing
finer-grained importance signals.
This would increase computational cost but allow a direct ablation of the
thermal-signal aspect of TCL.

\textbf{Class-incremental setting.}
Task-incremental evaluation assumes task identity at test time.
Removing this assumption requires a unified output head and changes the nature of
forgetting dramatically.
TCL's anchor strategy must be modified for this setting because task boundaries
in class-incremental training correspond to changes in the output space, not just
the input distribution.

\textbf{Online $D_{\mathrm{PR}}$ computation.}
A low-rank approximation of the gradient covariance would allow online $D_{\mathrm{PR}}$
estimation with bounded overhead, enabling the participation-ratio scaling that was
disabled here for efficiency.

\textbf{Longer task sequences.}
The five-task setting of Split-CIFAR-10 limits the number of anchor accumulations.
Benchmarks with 10--20 tasks would test whether penalty drift becomes a problem as
the number of anchored checkpoints grows.
